% !PRINTED 480
\weberpage{480}
We can now also generalise this theorem in the following way:

3. \gesperrt{If by the substitution $y = A(x)$ any form $\varphi(x)$ goes over into $\Phi(y)$, and if any second form $\psi(\xi)$ goes over by the transposed substitution $\xi = A_1(\eta)$ into $\Psi(\eta)$, then one obtains a new transformation by $A$ if one makes in $\psi(\xi)$ and $\Psi(\eta)$ the interchanges}

\[
\begin{aligned}
&\xi_1^\alpha \xi_2^\beta \xi_3^\gamma \quad \text{with} \quad \frac{\partial^{\alpha + \beta + \gamma} \varphi}{\partial x_1^\alpha \partial x_2^\beta \partial x_3^\gamma} \\
&\eta_1^\alpha \eta_2^\beta \eta_3^\gamma \quad \text{with} \quad \frac{\partial^{\alpha + \beta + \gamma} \Phi}{\partial y_1^\alpha \partial y_2^\beta \partial y_3^\gamma},
\end{aligned}
\]

\gesperrt{that is, if, as one may also express it, one puts in $\psi$ and $\Psi$ for the powers and products of the variables $\xi, \eta$ the corresponding derivatives of $\varphi, \Phi$ with respect to the variables $x, y$.}

From the rule of composition of linear substitutions there now follows immediately the following theorem:

4. \gesperrt{If $A$ runs through a group $G$, then the transposed substitution $A_1$ runs through a group $G_1$. If $A, B$ are two elements of $G$ and $A_1, B_1$ the corresponding elements of $G_1$, then $AB$ and $B_1 A_1$ are corresponding elements. The groups $G$ and $G_1$ are said to be contragredient to each other.}

The two groups $G$ and $G_1$ are, however, related to each other isomorphically only if one lets not the element $A$ but the element $A^{-1}$ correspond to $A_1$; for then $A_1 B_1$ corresponds to the element $A^{-1} B^{-1} = (BA)^{-1}$.

The invariants of the group $G_1$ are called contravariants of the group $G$. Accordingly the invariants of $G$ are also the contravariants of $G_1$.

From (3) there then follows the following theorem:

5. \gesperrt{If in a contravariant of $G$ one replaces the powers and products of the variables by the corresponding derivatives of an invariant, one again obtains an invariant of $G$}
